\section{Conclusion}

This work evaluated the performance scalability and adversarial robustness of Quantum Machine Learning models across several simulated noisy environments. Our results demonstrate that while density matrix simulation of noisy 8-qubit devices is computationally intensive, the implementation of multicore parallelization and the selection of optimized Aer solver methods make such studies feasible on modern classical hardware. We establish a quantitative scaling factor of ~10.86× for runtime increase when doubling the qubit count from 4 to 8.

Regarding adversarial security, we confirm that white-box perturbations generated via PGD can reduce classification accuracy from near-perfect levels to as low as 12\% in noisy environments. Our investigation into defense strategies reveals that standard adversarial retraining provides only modest recovery when using low ratios of adversarial samples. In contrast, a 50/50 adversarial-to-benign retraining split was shown to be highly effective, restoring robust accuracy to 84\%. We also demonstrate that Lipschitz gradient regularization provides a substantial 20\% absolute improvement in robustness over the standard retraining baseline. Furthermore, we find that architectural choices such as non-linear feature mappings can inadvertently increase vulnerability to gradient-based attacks. These findings emphasize the necessity of co-designing QNN architectures and training protocols to achieve practical security in the NISQ era.

\paragraph{Outlook}
Future research will explore the transferability of adversarial robustness across diverse quantum architectures. Specifically, we aim to investigate "cross-backend" resilience, where QNN models are trained on one QPU simulator and subsequently evaluated on different hardware-mimicking environments to assess the impact of varying noise profiles. Additionally, we plan to integrate advanced model optimization techniques, such as Approximate State Preparation (ASP) \cite{West2023}, to further enhance the defensive capabilities of QNNs against sophisticated adversarial perturbations.
